\chapter{Construction de la courbe}\label{ch:rc:curve-construction}

Un vendredi soir, un desk de taux remplace l'interpolation de sa courbe en
dollars~: les taux zéro étaient reliés par des droites, et à partir de lundi ils seront reliés
par la méthode que le Trésor américain lui-même a adoptée en 2021. Chaque
cotation d'entrée du lundi matin est exactement là où elle était le vendredi,
chaque instrument d'entrée se reprice toujours au dernier décimal, et pourtant le carnet
de sept swaps hors \omterm{def:rc:curve-construction:pillar}{pilier} du desk, pour cinq milliards de dollars de notionnel, affiche une perte de
90 mille dollars, et la couverture que réclame le système de risque sur le swap à quinze ans
a augmenté de deux tiers. Rien sur le marché n'a bougé~; c'est la courbe entre
les cotations, qui a bougé. Une courbe est un modèle, ajusté à une poignée de prix et utilisé pour
valoriser tout ce qui se trouve entre eux. Ce chapitre la construit à partir des instruments
de bout en bout, et mesure le coût de chaque choix en valeur et en couvertures. Le
livre~2, chapitre~9, a bootstrapé une seule courbe de swaps au jour le jour, avec
une seule interpolation~; ici la courbe prend ensemble les dépôts, les contrats à terme et les swaps, est
résolue globalement, et interpolée de quatre façons.

\section{Choix des instruments}

\begin{definition}[Pilier, taux forward instantané]\label{def:rc:curve-construction:pillar}
Un \emph{pilier}\index{pilier} de courbe est une date à laquelle la courbe porte un
paramètre libre, ici le taux zéro à capitalisation continue $z(T)$, de sorte que
$P(T)=e^{-z(T)T}$~; la courbe entre deux piliers est fixée par une règle d'interpolation.
Le \emph{taux forward instantané}\index{taux forward instantané} observé en
$t$ pour l'échéance $T$ est
\[
f(t,T) = -\frac{\partial}{\partial T}\ln P(t,T),
\qquad
P(t,T)=\exp\Bigl(-\int_t^T f(t,u)\,du\Bigr),
\]
le taux auquel on peut figer, en $t$, un emprunt sur $[T, T+dT]$.
\end{definition}

\begin{definition}[Calibration de la courbe]\label{def:rc:curve-construction:calibration}
La \emph{calibration de la courbe}\index{calibration de la courbe} est la solution pour les paramètres
des \omterm{def:rc:curve-construction:pillar}{piliers} de la courbe, telle qu'un jeu choisi d'instruments, chacun
évalué sur la courbe, reproduise sa cotation de marché. Avec un \omterm{def:rc:curve-construction:pillar}{pilier} par
instrument, c'est un système carré d'équations~; la notion générale de
calibration, qui consiste à ajuster les paramètres d'un modèle aux prix, est traitée dans le
livre~4, chapitre~24.
\end{definition}

Les instruments sont choisis pour leur liquidité et de façon que chaque échéance soit couverte
une seule fois. En tête d'une courbe en dollars se trouvent les fixings au jour le jour et les
dépôts courts au jour le jour~; puis les contrats à terme SOFR à trois mois, dont les
trimestres de référence couvrent la première ou les deux premières années~; puis des swaps OIS au pair
(livre~2, chapitres~8 et~9). La règle est \emph{sans recouvrement}~: si quatre
contrats à terme fixent déjà les forwards jusqu'à décembre 2027, un swap d'un an échéant en
septembre 2027 ajoute un second prix, légèrement différent, pour une période que la courbe
a déjà valorisée. Deux cotations pour une même période ne peuvent être toutes deux reproduites, sauf
si elles concordent~; le constructeur retire l'une des deux ou ajuste les deux au sens des moindres carrés
et renonce à reproduire exactement les prix. La \Cref{fig:rc:curve-construction:instruments}
montre le choix du chapitre.

\begin{omfigure}
\begin{tikzpicture}[x=4.4cm,y=1cm,font=\small,
  bar/.style={draw=omDef,fill=omDef!15,thick,minimum height=0.36cm},
  lab/.style={font=\footnotesize,anchor=west}]
\draw[->,thick] (0,-1.2) -- (2.25,-1.2) node[right] {années};
\foreach \x/\l in {0/0,0.5/0.5,1/1,1.5/1.5,2/2} {\draw (\x,-1.28) -- (\x,-1.12); \node[below,font=\footnotesize] at (\x,-1.28) {\l};}
% deposits
\draw[bar] (0,2.55) rectangle (0.082,2.85); \node[lab] at (0.10,2.70) {dépôt 1M};
\draw[bar] (0,2.05) rectangle (0.249,2.35); \node[lab] at (0.27,2.20) {dépôt 3M};
% futures
\foreach \a/\b/\n in {0.2137/0.463/F1,0.463/0.7123/F2,0.7123/0.9616/F3,0.9616/1.211/F4}
  {\draw[draw=omThm,fill=omThm!15,thick] (\a,1.05) rectangle (\b,1.45); \node[font=\footnotesize] at ({(\a+\b)/2},1.25) {\n};}
\node[lab] at (1.24,1.25) {contrats à terme SOFR (trimestres IMM)};
% swaps
\draw[draw=omProp,fill=omProp!15,thick] (0,0.05) rectangle (2.0027,0.45); \node[font=\footnotesize] at (1.0,0.25) {swap OIS au pair 2A};
\node[lab,align=left] at (2.03,0.25) {puis 3A, 4A, 5A, 7A,\\10A, \dots, 30A};
\draw[dashed,gray,thick] (0,-0.75) rectangle (1.0,-0.45); \node[lab,gray] at (1.03,-0.6) {swap 1A~: non utilisé, recouvre F1--F4};
% pillars
\foreach \x in {0.082,0.249,0.463,0.7123,0.9616,1.211,2.0027} {\fill[omExo] (\x,-1.2) circle (1.6pt);}
\end{tikzpicture}

\omcaption{Les instruments de la courbe en dollars du chapitre sur ses deux premières
années, et leurs \omterm{def:rc:curve-construction:pillar}{piliers} (points sur l'axe, un à l'échéance de chaque
instrument). Les contrats à terme commencent le troisième mercredi de décembre 2026 et couvrent
cinq trimestres~; le swap d'un an porterait sur une période que les contrats à terme
valorisent déjà, et il est laissé de côté.}\label{fig:rc:curve-construction:instruments}
\end{omfigure}

\begin{dated}{2026-09}{Courbes publiées}\label{dat:rc:curve-construction:published}
La courbe quotidienne officielle des taux au pair du Trésor américain est construite avec une
\omterm{def:rc:curve-construction:monotone}{interpolation monotone convexe} depuis le 6 décembre 2021, en remplacement d'une spline de
Hermite quasi-cubique~; ses entrées sont les cotations indicatives au cours acheteur des
bons, billets et obligations adjugés le plus récemment, recueillies par la New York Fed vers 15:30.
La BCE publie deux courbes gouvernementales de la zone euro, à partir des obligations notées AAA et
de toutes les obligations des administrations centrales, chaque jour ouvré TARGET à 12:00 CET,
chacune sous la forme de six paramètres d'un modèle de Svensson
(\cref{def:rc:curve-construction:ns}) et les taux qu'elles impliquent.
\end{dated}

\section{Calibration~: bootstrap de la courbe et ajustement global}

La cotation de chaque instrument est une fonction des facteurs d'actualisation jusqu'à son
échéance. Un dépôt du spot à $T$ au taux simple $R$ (ACT/360) suppose
$P(T) = 1/(1+R\delta)$~; un contrat à terme sur le trimestre $[T_1,T_2]$, une fois
son ajustement de convexité retiré (livre~2, chapitre~8), fixe
$P(T_1)/P(T_2) = 1+F\delta$~; un swap au pair fixe $P(T_0)-P(T_n)=K\,\sum_i
\delta_iP(t_i)$. Avec $n$ instruments et $n$ \omterm{def:rc:curve-construction:pillar}{piliers}, le système
$m_i(z_1,\dots,z_n)=q_i$ est carré.

\begin{proposition}[Quand le bootstrap de la courbe fonctionne]\label{prop:rc:curve-construction:triangular}
Si l'interpolation est \emph{locale}, au sens où la courbe sur
$[T_{j-1},T_j]$ ne dépend que des \omterm{def:rc:curve-construction:pillar}{piliers} $j-1$ et $j$, et chaque instrument
$i$ a pour échéance $T_i$, alors $\partial m_i/\partial z_j = 0$ pour $j>i$~: la
jacobienne est triangulaire inférieure, et le système peut être résolu un \omterm{def:rc:curve-construction:pillar}{pilier} à la
fois, chacune étant une recherche de racine à une dimension. C'est le bootstrap de la courbe. Les interpolations
linéaires des taux zéro et des logarithmes des facteurs d'actualisation sont locales~; une
spline cubique ne l'est pas, et la méthode monotone convexe est locale sauf par les
forwards qu'elle estime aux deux \omterm{def:rc:curve-construction:pillar}{piliers} voisins.
\end{proposition}
\begin{proof}
L'instrument $i$ ne lit les facteurs d'actualisation qu'aux dates jusqu'à $T_i$~; sous une règle
locale, ceux-ci ne dépendent que des \omterm{def:rc:curve-construction:pillar}{piliers} $1,\dots,i$.
\end{proof}

\begin{method}[Calibration globale de la courbe par la méthode de Newton]\label{met:rc:curve-construction:newton}
Partir de $z^{(0)}_j=q_j$. À chaque itération, calculer les résidus $r_i = m_i(z)
- q_i$ et la jacobienne $J_{ij}=\partial m_i/\partial z_j$ par
différences finies~; mettre à jour $z \leftarrow z - J^{-1}r$~; s'arrêter lorsque $\max_i|r_i|<10^{-12}$.
Le même code sert à toutes les interpolations, locales ou non, et la jacobienne
finale est conservée~: elle convertit les sensibilités aux taux des \omterm{def:rc:curve-construction:pillar}{piliers} en
sensibilités aux cotations (\cref{ch:rc:rates-risk}).
\end{method}

\omcode{code/firm/curvebuild/firm_curvebuild.py}{256}{284}{Calibration globale de tous les taux zéro des piliers par la méthode de Newton.}\label{lst:rc:curve-construction:calibrate}

Sur les dix-sept instruments de la \cref{fig:rc:curve-construction:instruments},
la méthode de Newton converge en trois itérations pour les quatre interpolations,
avec des résidus inférieurs à $10^{-15}$. La vitesse n'est pas l'enjeu ici~; ce
que permet la forme globale, c'est de changer l'interpolation sans changer le solveur.

\section{Interpolation et courbe à terme}

Chaque interpolation reprice les données d'entrée~; elles diffèrent par ce qu'elles
disent entre les \omterm{def:rc:curve-construction:pillar}{piliers}, et c'est dans la courbe à terme que les écarts apparaissent.

\begin{proposition}[Forwards à partir des taux zéro]\label{prop:rc:curve-construction:fwd}
Avec $y(T) = z(T)T = -\ln P(T)$, le \omterm{def:rc:curve-construction:pillar}{taux forward instantané} vaut $f(0,T) =
y'(T) = z(T) + T z'(T)$. Le forward moyen sur un intervalle entre deux \omterm{def:rc:curve-construction:pillar}{piliers} est fixé par
les deux seuls \omterm{def:rc:curve-construction:pillar}{piliers}~: $\bar f_j = (y_j - y_{j-1})/(T_j-T_{j-1})$, le
\emph{forward discret}, sans index ici.
\end{proposition}

\begin{definition}[Interpolation à forward constant]\label{def:rc:curve-construction:flatfwd}
L'\emph{interpolation à forward constant}\index{interpolation à forward constant} rend
$\ln P(T)$ linéaire entre les \omterm{def:rc:curve-construction:pillar}{piliers}~: le \omterm{def:rc:curve-construction:pillar}{taux forward instantané} est constant sur
chaque intervalle et égal à son forward discret, et il saute aux \omterm{def:rc:curve-construction:pillar}{piliers}.
\end{definition}

L'interpolation linéaire des taux zéro paraît plus lisse sur un graphique des taux zéro et
elle est moins bonne sur un graphique des forwards~: d'après la \cref{prop:rc:curve-construction:fwd},
le forward vaut $z + Tz'$, et $z'$ saute à chaque \omterm{def:rc:curve-construction:pillar}{pilier}, multiplié par $T$. Au
\omterm{def:rc:curve-construction:pillar}{pilier} de deux ans du chapitre ($z = 3.2883\%$), la courbe zéro descend à sa gauche
et remonte à sa droite~; le forward passe de $2.964~\%$ à $3.331~\%$, un
saut de 37~points de base produit par la règle et non par le marché. Une spline
cubique naturelle sur les taux zéro (livre~4, chapitre~28) supprime les
sauts mais, étant un ajustement global, laisse une modification à un \omterm{def:rc:curve-construction:pillar}{pilier} se
propager à tous les autres.

\begin{definition}[Interpolation monotone convexe]\label{def:rc:curve-construction:monotone}
L'\emph{interpolation monotone convexe}\index{monotone convex interpolation}
(Hagan et West, 2006) construit directement la courbe à terme. Sur chaque intervalle,
elle conserve le forward discret $\bar f_j$, estime le \omterm{def:rc:curve-construction:pillar}{taux forward instantané}
$f_j$ à chaque \omterm{def:rc:curve-construction:pillar}{pilier} comme la moyenne pondérée par le temps des deux forwards discrets
adjacents, et remplit l'intervalle avec $f(T)=\bar f_j + G(x)$,
$x=(T-T_{j-1})/(T_j-T_{j-1})$, où $G$ est une fonction quadratique, ou une quadratique
raccordée à une constante, avec $G(0)=f_{j-1}-\bar f_j$, $G(1)=f_j-\bar f_j$ et
$\int_0^1G=0$, choisie parmi quatre cas de façon que $G$ ne dépasse jamais ses valeurs
aux extrémités. La courbe à terme est continue, reproduit chaque forward discret
et reste dans l'intervalle des forwards voisins.
\end{definition}

\omcode{code/firm/curvebuild/firm_curvebuild.py}{123}{138}{Les quatre cas de l'interpolant monotone convexe sur un intervalle.}\label{lst:rc:curve-construction:mc}

\begin{omfigure}
\begin{tikzpicture}
\begin{axis}[width=12cm,height=6cm,
  xlabel={maturité $T$ (années)},ylabel={$f(0,T)$ (\%)},
  tick label style={font=\small},label style={font=\small},
  xmin=0,xmax=30,ymin=2.9,ymax=4.75,grid=major,grid style={gray!25},
  legend columns=4,legend style={font=\footnotesize,at={(0.5,-0.25)},anchor=north,draw=none,fill=none,/tikz/every even column/.append style={column sep=8pt}},
  table/col sep=comma]
\addplot[gray,thick] table[x=t,y=lz]{figdata/rates-credit-risk/01-curve-construction/forwards.csv};
\addplot[omDef,thick] table[x=t,y=ff]{figdata/rates-credit-risk/01-curve-construction/forwards.csv};
\addplot[omProp,thick,dashed] table[x=t,y=cs]{figdata/rates-credit-risk/01-curve-construction/forwards.csv};
\addplot[omThm,very thick] table[x=t,y=mc]{figdata/rates-credit-risk/01-curve-construction/forwards.csv};
\legend{zéro linéaire,forward constant,cubique sur les zéros,monotone convexe}
\end{axis}
\end{tikzpicture}

\omcaption{Quatre courbes de \omterm{def:rc:curve-construction:pillar}{taux forward instantanés} issues des mêmes dix-sept cotations~;
chacune reprice exactement chaque donnée. Les taux zéro linéaires donnent un profil en dents
de scie~; les forwards constants donnent des marches~; la spline cubique et la courbe monotone
convexe sont continues et quasiment confondues à cette échelle~: elles diffèrent par leur réaction
à la modification d'une seule cotation (\cref{fig:rc:curve-construction:locality}).
Données~: la courbe en dollars illustrative du chapitre et
tutoriel.}\label{fig:rc:curve-construction:forwards}
\end{omfigure}

Pour les obligations d'État, nombreuses, bruitées et pas toutes au pair, on utilise
plutôt une courbe paramétrique lisse ajustée par moindres carrés, et non un \omterm{def:rc:curve-construction:pillar}{pilier} par
instrument.

\begin{definition}[Courbes de Nelson--Siegel et de Svensson]\label{def:rc:curve-construction:ns}
Une \emph{courbe de Nelson--Siegel}\index{courbe de Nelson--Siegel} est la courbe de taux zéro
\[
z(T) = \beta_0 + \beta_1\,\frac{1-e^{-T/\tau_1}}{T/\tau_1}
 + \beta_2\Bigl(\frac{1-e^{-T/\tau_1}}{T/\tau_1}-e^{-T/\tau_1}\Bigr),
\]
avec $\beta_0$ pour niveau de long terme, $\beta_0+\beta_1$ pour extrémité courte et une bosse d'amplitude
$\beta_2$ située près de $\tau_1$~; sa courbe forward est $f(0,T) = \beta_0 + \beta_1
e^{-T/\tau_1} + \beta_2 (T/\tau_1)e^{-T/\tau_1}$. L'extension de Svensson ajoute un
second terme de bosse $\beta_3$ avec son propre $\tau_2$.
\end{definition}

\begin{example}[La courbe AAA de la BCE]\label{ex:rc:curve-construction:ecb}
Le 22 septembre 2026, les paramètres de Svensson de la courbe AAA de la BCE étaient $\beta_0=1.6306$,
$\beta_1=0.6958$, $\beta_2=2.5849$, $\beta_3=6.6622$ (en pour cent), $\tau_1=1.0632$,
$\tau_2=14.4456$ années. Ils donnent un taux zéro à dix ans de $3.4527~\%$, qui est
le taux spot à dix ans publié par la BCE, ainsi que les taux à un, cinq et trente ans
publiés à six décimales~; le \omterm{def:rc:curve-construction:pillar}{taux forward instantané} monte de
$2.33~\%$ à l'extrémité courte jusqu'à un sommet de $4.08~\%$ vers quatorze ans
(\cref{fig:rc:curve-construction:svensson}).
\end{example}

\begin{omfigure}
\begin{tikzpicture}
\begin{axis}[width=11cm,height=5.4cm,
  xlabel={échéance $T$ (années)},ylabel={taux (\%)},
  tick label style={font=\small},label style={font=\small},
  xmin=0,xmax=30,ymin=2,ymax=4.4,grid=major,grid style={gray!25},
  legend columns=3,legend style={font=\footnotesize,at={(0.5,-0.25)},anchor=north,draw=none,fill=none,/tikz/every even column/.append style={column sep=8pt}},
  table/col sep=comma]
\addplot[omDef,very thick] table[x=t,y=zero]{figdata/rates-credit-risk/01-curve-construction/svensson.csv};
\addplot[omThm,thick,dashed] table[x=t,y=fwd]{figdata/rates-credit-risk/01-curve-construction/svensson.csv};
\addplot[only marks,mark=*,mark size=2.2pt,omExo] table[x=t,y=zero]{figdata/rates-credit-risk/01-curve-construction/ecbpoints.csv};
\legend{taux zéro,forward instantané,taux spot publiés par la BCE}
\end{axis}
\end{tikzpicture}

\omcaption{La courbe souveraine AAA de la zone euro du 22 septembre 2026, calculée à partir des
six paramètres de Svensson que publie la BCE~; les points sont les taux spot que la BCE
publie pour un, cinq, dix et trente ans, que la formule reproduit.
Source~: statistiques de la BCE (ECB Data Portal, série
YC.B.U2.EUR.4F.G\_N\_A.SV\_C\_YM).}\label{fig:rc:curve-construction:svensson}
\end{omfigure}

\section{La partie courte~: contrats à terme, passages de fin de période et dates de réunion}

Les deux premières années d'une courbe en dollars sont construites à partir d'instruments qui valorisent
des périodes à terme, et non des taux au comptant. Un contrat à terme sur un trimestre IMM fixe le
taux forward d'un trimestre, après son ajustement de convexité~; quatre d'entre eux fixent quatre
forwards discrets consécutifs, et l'interpolation ne détermine que la forme à l'intérieur de chaque
trimestre. Deux caractéristiques de la partie courte ne sont pas lisses et ne doivent pas être
interpolées comme si elles l'étaient. Le passage
de fin de période du livre~2, chapitre~2, est un saut du taux au jour le jour lors d'une fin d'année
ou de trimestre~: un constructeur de courbe qui en tient compte ajoute un \omterm{def:rc:curve-construction:pillar}{pilier} supplémentaire
autour de la date de passage, ou un spread distinct sur ces jours, afin que cette nuit coûteuse ne soit
pas diluée sur tout un trimestre. La trajectoire implicite des taux directeurs
est une fonction en escalier~: le taux au jour le jour change lors des réunions des banques centrales et non
entre elles, si bien que les courbes de la partie courte sont souvent construites avec une
\omterm{def:rc:curve-construction:flatfwd}{interpolation à forward constant} entre les dates de réunion, avec un \omterm{def:rc:curve-construction:pillar}{pilier} par réunion (construction
du livre~2, \texttt{firm.meetings}). L'interpolation lisse convient là où le marché n'a pas de
calendrier, à partir de deux ou trois ans.

\begin{remark}[Quelle courbe pour quel usage]\label{rem:rc:curve-construction:which}
Les forwards constants entre les dates de réunion à court terme, une règle lisse telle que la
monotone convexe au-delà, et une courbe paramétrique pour la valeur relative dans les
obligations d'État (une obligation chère ou bon marché par rapport à la courbe ajustée) sont des choix complémentaires,
et non concurrents. L'un quelconque d'entre eux échoue si un instrument est inclus deux fois pour une
même période, ou si une cotation est périmée.
\end{remark}

\section{Conséquences de l'interpolation sur la couverture}

\begin{definition}[Localité de l'interpolation]\label{def:rc:curve-construction:locality}
Un schéma d'interpolation possède la \emph{localité de l'interpolation}\index{localité de l'interpolation}
lorsqu'une variation d'une cotation d'entrée ne modifie la courbe qu'au voisinage de l'échéance
de cette entrée. La localité détermine l'allure d'une couverture par buckets~: avec un schéma local, un
swap hors \omterm{def:rc:curve-construction:pillar}{pilier} est couvert par les \omterm{def:rc:curve-construction:pillar}{piliers} qui l'entourent~; avec un schéma non local, il acquiert
un risque, de signe quelconque, sur des \omterm{def:rc:curve-construction:pillar}{piliers} éloignés
de son échéance.
\end{definition}

\begin{omfigure}
\begin{tikzpicture}
\begin{axis}[width=12cm,height=5.8cm,
  xlabel={échéance $T$ (années)},ylabel={variation de $f(0,T)$ (pb)},
  tick label style={font=\small},label style={font=\small},
  xmin=0,xmax=20,ymin=-5.5,ymax=5.5,grid=major,grid style={gray!25},
  legend columns=4,legend style={font=\footnotesize,at={(0.5,-0.25)},anchor=north,draw=none,fill=none,/tikz/every even column/.append style={column sep=8pt}},
  table/col sep=comma]
\addplot[gray,thick] table[x=t,y=lz]{figdata/rates-credit-risk/01-curve-construction/locality.csv};
\addplot[omDef,thick] table[x=t,y=ff]{figdata/rates-credit-risk/01-curve-construction/locality.csv};
\addplot[omProp,thick,dashed] table[x=t,y=cs]{figdata/rates-credit-risk/01-curve-construction/locality.csv};
\addplot[omThm,very thick] table[x=t,y=mc]{figdata/rates-credit-risk/01-curve-construction/locality.csv};
\legend{zéro linéaire,forward constant,cubique sur les zéros,monotone convexe}
\end{axis}
\end{tikzpicture}

\omcaption{Réponse de la courbe forward à une hausse d'un point de base de la cotation du swap à
sept ans, la courbe étant recalibrée. Les forwards constants ne bougent qu'entre cinq et dix ans
(à la hausse avant sept ans, à la baisse après, de sorte que le swap à dix ans se valorise toujours
au prix de marché)~; la courbe monotone convexe ajoute de petites
ondulations juste en dehors de cet intervalle~; la spline cubique fait bouger les forwards de trois
ans jusqu'au-delà de vingt. Données~: le tutoriel du chapitre.}\label{fig:rc:curve-construction:locality}
\end{omfigure}

\begin{example}[Un swap de huit ans sur quatre courbes]\label{ex:rc:curve-construction:eight}
Un swap payeur au pair de 100~millions de USD sur huit ans a un DV01 parallèle de
d'environ 69\,600 USD sur les quatre courbes~: elles s'accordent sur le niveau. Elles diffèrent
sur leur localisation (\cref{fig:rc:curve-construction:buckets}). Les forwards constants placent
40\,400 USD sur la cotation à sept ans et 29\,200 USD sur celle à dix ans~; la spline cubique
place 62\,100 USD sur sept ans, $-22\,200$ sur cinq, $+11\,000$ sur quatre
et $-9\,500$ sur douze~; la monotone convexe se situe entre les deux. Chaque ensemble de buckets
somme au DV01 parallèle, et chacun appelle une couverture différente.
\end{example}

\begin{omfigure}
\begin{tikzpicture}
\begin{axis}[width=12cm,height=5.4cm,ybar,bar width=7pt,
  ylabel={DV01 (milliers de USD par pb)},
  tick label style={font=\small},label style={font=\small},
  xmin=-0.6,xmax=6.6,ymin=-30,ymax=70,grid=major,grid style={gray!25},
  xtick=data,xticklabels={3A,4A,5A,7A,10A,12A,15A},
  legend columns=3,legend style={font=\footnotesize,at={(0.5,-0.2)},anchor=north,draw=none,fill=none,/tikz/every even column/.append style={column sep=8pt}},
  table/col sep=comma]
\addplot[fill=omDef!55,draw=omDef] table[x=i,y=ff]{figdata/rates-credit-risk/01-curve-construction/buckets8.csv};
\addplot[fill=omProp!55,draw=omProp] table[x=i,y=cs]{figdata/rates-credit-risk/01-curve-construction/buckets8.csv};
\addplot[fill=omThm!55,draw=omThm] table[x=i,y=mc]{figdata/rates-credit-risk/01-curve-construction/buckets8.csv};
\legend{forward constant,cubique sur les zéros,monotone convexe}
\end{axis}
\end{tikzpicture}

\omcaption{DV01 par buckets d'un swap payeur au pair de huit ans de 100~millions de USD, par cotation
d'entrée, pour trois interpolations (chaque cotation étant décalée d'un centième de point de base, courbe
recalibrée, résultat ramené à un point de base). La date de huit ans se situe entre les \omterm{def:rc:curve-construction:pillar}{piliers} de sept et
dix ans~; les forwards constants ne chargent que ces deux-là, la monotone convexe ajoute de petits buckets
à cinq et douze ans, la spline cubique de grands buckets des deux signes
de trois à quinze ans. Données~: le tutoriel
du chapitre.}\label{fig:rc:curve-construction:buckets}
\end{omfigure}

\begin{remark}[La monotone convexe n'est pas linéaire en ses entrées]\label{rem:rc:curve-construction:nonlinear}
Les courbes à taux zéro linéaires, à forward constant et spline cubique sont des fonctions linéaires des
taux des \omterm{def:rc:curve-construction:pillar}{piliers}, donc un DV01 par buckets ne dépend pas de la taille du choc. La
courbe monotone convexe ne l'est pas~: lequel de ses quatre cas s'applique sur un intervalle
dépend des entrées, si bien qu'un choc d'un point de base peut faire changer de cas et modifier
la forme du forward à l'intérieur d'un intervalle d'un montant supérieur au choc lui-même. La valeur
reste continue, mais les buckets varient avec la taille du choc. Le chapitre
applique donc un choc d'un centième de point de base, puis remet le résultat à l'échelle d'un point de base~; l'exercice~7
mesure cet écart.
\end{remark}

\section{Tutoriel~: un seul jeu de cotations, quatre courbes}\label{tut:rc:curve-construction:tutorial}

\begin{tutorial}
\textbf{Objectif.} Calibrer une courbe en dollars à partir de dépôts, de contrats à terme et de swaps
selon quatre interpolations, et comparer les forwards, la \omterm{def:rc:curve-construction:locality}{localité de l'interpolation} et les couvertures.
\textbf{État final~:}
\cref{fig:rc:curve-construction:forwards,fig:rc:curve-construction:locality,fig:rc:curve-construction:buckets}
et les chiffres de l'\cref{ex:rc:curve-construction:eight}.
\begin{tutsteps}
\item \textbf{Instruments.} Deux dépôts, quatre contrats à terme, du troisième
mercredi de décembre 2026, onze swaps~; aucun swap à un an.
\omcode{code/rates-credit-risk/01-curve-construction/python/rc_curves.py}{30}{36}{L'ensemble des instruments de la courbe du chapitre.}\label{lst:rc:curve-construction:instruments}
\item \textbf{Calibrer} avec \texttt{calibrate(SPOT, instruments(), kind)} pour chacun
des quatre types~; vérifier que le résidu maximal est inférieur à $10^{-12}$ et que le
nombre d'itérations vaut trois.
\item \textbf{Forwards.} Tabuler \texttt{curve.fwd\_t(t)} tous les 0.05~an~; on
doit voir la dent de scie des taux zéro linéaires au \omterm{def:rc:curve-construction:pillar}{pilier} à deux ans (2.964~\% puis
3.331~\%).
\item \textbf{Localité et buckets.} Lancer \texttt{locality} pour le \omterm{def:rc:curve-construction:pillar}{pilier} \texttt{7Y} et
\texttt{off\_pillar\_buckets(8)}~; \texttt{fig\_rc\_curves.py} écrit les graphiques.
\end{tutsteps}
\textbf{Pour aller plus loin.} Remettre le swap à un an, à deux points de base
des contrats à terme, et regarder la courbe à terme de septembre 2027~; appliquer un choc d'un point de base entier
sous \omterm{def:rc:curve-construction:monotone}{interpolation monotone convexe} et comparer les buckets à ceux de
la \cref{fig:rc:curve-construction:buckets}.
\end{tutorial}

\section{Construction~: un constructeur de courbe multi-instruments}\label{bld:rc:curve-construction:curvebuild}

\begin{build}
\bfield{Rôle} La première courbe de la bibliothèque du livre~6~: tout instrument de taux,
toute simulation d'exposition et tout calcul de risque du livre lit une courbe construite ici, et la
bibliothèque de valorisation du livre~5, chapitre~28, l'accepte comme
courbe d'actualisation.
\bfield{Interface} \texttt{Deposit}, \texttt{Future}, \texttt{Swap} avec
\texttt{quote()} et \texttt{model(curve)}~; \texttt{calibrate(spot, instruments,
kind)} renvoyant la courbe, le jacobien et les diagnostics~; \texttt{ZeroCurve}
avec \texttt{df}, \texttt{zero\_t}, \texttt{fwd\_t} et \texttt{bumped(pillar,
size)}~; \texttt{bucket\_sensitivities}~; types \texttt{linear\_zero},
\texttt{flat\_forward}, \texttt{cubic\_zero}, \texttt{monotone\_convex}.
\bfield{Règles} \omterm{def:rc:curve-construction:pillar}{Piliers} aux échéances des instruments, strictement croissants (un
chevauchement est une erreur)~; les temps en ACT/365, les accruals en ACT/360~; forward constant au-delà du
dernier \omterm{def:rc:curve-construction:pillar}{pilier}~; \texttt{firm\_curve} du livre~2 importée, non modifiée.
\bfield{Tests d'acceptation} \texttt{code/firm/curvebuild/tests/}~: chaque
instrument se revalorise sous tous les types~; l'\omterm{def:rc:curve-construction:flatfwd}{interpolation à forward constant} sur les swaps égale le
bootstrap du livre~2~; des cotations au pair plates donnent des forwards constants~; le forward monotone convexe est
continu aux \omterm{def:rc:curve-construction:pillar}{piliers}~; les types locaux ont un jacobien triangulaire inférieur, contrairement à la
spline cubique~; un choc à sept ans se propage au-delà de quinze ans seulement sous la
spline~; \texttt{bumped} décale un \omterm{def:rc:curve-construction:pillar}{pilier} ou tous.
\bfield{Pour aller plus loin} \omterm{def:rc:curve-construction:pillar}{Piliers} de passage de fin de période et \omterm{def:rc:curve-construction:pillar}{piliers} aux dates de réunion~; un mode aux moindres carrés avec plus
d'instruments que de \omterm{def:rc:curve-construction:pillar}{piliers}~; l'amendement de positivité de la
méthode monotone convexe~; un jacobien analytique.
\end{build}

\begin{omsources}
\item P.\ S.\ Hagan and G.\ West, ``Interpolation methods for curve
construction'', \emph{Applied Mathematical Finance} 13(2), 2006.
\item C.\ R.\ Nelson and A.\ F.\ Siegel, ``Parsimonious modeling of yield
curves'', \emph{Journal of Business} 60(4), 1987; L.\ E.\ O.\ Svensson,
``Estimating and interpreting forward interest rates: Sweden 1992--1994'', NBER
Working Paper 4871, 1994.
\item US Department of the Treasury, \emph{Treasury Yield Curve Methodology};
ECB, \emph{Euro area yield curves} (ECB Data Portal).
\item L.\ Andersen and V.\ Piterbarg, \emph{Interest Rate Modeling}, tome 1,
Atlantic Financial Press, 2010, chapitre~6.
\end{omsources}

\section{Exercices}

\begin{exercise}[$\star$]\label{exo:rc:curve-construction:1}
Les taux zéro à capitalisation continue à deux et trois ans sont $3.25~\%$
et $3.30\%$. Sous \omterm{def:rc:curve-construction:flatfwd}{interpolation à forward constant}, quel est le
\omterm{def:rc:curve-construction:pillar}{taux forward instantané} à $2.5$ ans~?
\end{exercise}

\begin{exercise}[$\star$]\label{exo:rc:curve-construction:2}
Un contrat à terme SOFR sur le trimestre de juin à septembre 2027 se négocie à $96.72$ et
son ajustement de convexité vaut $0.6$ points de base. Quel taux forward la courbe doit
reproduire pour ce trimestre~?
\end{exercise}

\begin{exercise}[$\star$]\label{exo:rc:curve-construction:3}
Pourquoi la courbe du chapitre laisse-t-elle de côté le swap à un an, et que
se passerait-il pour la calibration s'il était remis avec une cotation éloignée de deux points de base
de ce qu'impliquent les contrats à terme~?
\end{exercise}

\begin{exercise}[$\star\star$]\label{exo:rc:curve-construction:4}
Sur la courbe à taux zéro linéaires du chapitre, les \omterm{def:rc:curve-construction:pillar}{piliers} autour de deux ans sont situés en
$(1.2110, 3.4165~\%)$, $(2.0027, 3.2883~\%)$ et $(3.0027, 3.3095~\%)$ (années,
taux zéro). Calculer le \omterm{def:rc:curve-construction:pillar}{taux forward instantané} juste avant et juste après le
\omterm{def:rc:curve-construction:pillar}{pilier} à deux ans.
\end{exercise}

\begin{exercise}[$\star\star$]\label{exo:rc:curve-construction:5}
À partir des paramètres de la BCE de l'\cref{ex:rc:curve-construction:ecb}, calculer le
\omterm{def:rc:curve-construction:pillar}{taux forward instantané} à dix ans et la limite du taux zéro lorsque $T\to 0$.
\end{exercise}

\begin{exercise}[$\star\star$]\label{exo:rc:curve-construction:6}
Lire la \cref{fig:rc:curve-construction:buckets}. Un trader couvre le swap à huit ans
avec les buckets de la spline, puis le desk passe à l'\omterm{def:rc:curve-construction:flatfwd}{interpolation à forward constant}. Quelles
couvertures deviennent inutiles, et pourquoi la spline les réclamait-elle~?
\end{exercise}

\begin{exercise}[$\star\star\star$]\label{exo:rc:curve-construction:7}
\emph{Programmation.} Avec \texttt{off\_pillar\_buckets(8, bump=1e-4)}, calculer les
buckets du swap à huit ans, en \omterm{def:rc:curve-construction:monotone}{interpolation monotone convexe}, avec un bump d'un point de base entier
et leur somme~; comparer avec la somme pour un bump d'un centième de point de base et avec le
DV01 parallèle, puis expliquer.
\end{exercise}

\begin{exercise}[$\star\star\star$]\label{exo:rc:curve-construction:8}
\emph{Trouver la faille.} «~Notre nouvelle courbe réévalue chaque instrument d'entrée à $10^{-12}$ près, donc
ses prix pour tout autre instrument sont justes.~» Corriger cette affirmation.
\end{exercise}

\section{Problème~: La perte du lundi}

\begin{problem}[{Problème du week-end --- une interpolation modifiée pendant le week-end}]\label{pb:rc:curve-construction:1}
Le book d'un desk détient sept swaps au pair négociés sur la courbe du chapitre sous
interpolation linéaire des taux zéro~: payeurs de 6, 9, 13 et 22 ans (0.9, 0.6, 0.5
et 0.8 milliard de dollars) et receveurs de 8, 11 et 17 ans (0.7, 0.8 et 0.7
milliard de dollars), soit 5.0 milliards de dollars au total. Pendant le week-end, la courbe est reconstruite à partir des
mêmes cotations avec une \omterm{def:rc:curve-construction:monotone}{interpolation monotone convexe}.

\medskip\noindent\textbf{Partie I --- Le changement.}
\begin{enumerate}
\item Pourquoi chaque swap vaut-il exactement zéro le vendredi~?
\item Un instrument d'entrée change-t-il de valeur le lundi~? Pourquoi pas~?
\item Le taux au pair à 17 ans passe de $3.9975~\%$ à $4.0074~\%$. Que révèle cela
sur les forwards entre 15 et 20 ans~?
\item Quels swaps du book sont exposés au changement, et lesquels ne le sont pas~?
\item Donner le taux au pair à six ans avant et après, et expliquer pourquoi il
varie à peine.
\end{enumerate}

\medskip\noindent\textbf{Partie II --- Le P\&L.}
\begin{enumerate}[resume]
\item Donner la valeur du book le lundi.
\item S'agit-il d'un gain ou d'une perte, et quels swaps en produisent l'essentiel~?
\item Le desk doit-il le comptabiliser en P\&L de trading~? Qui doit trancher~?
\item Comment le présenter au responsable du desk en une phrase~?
\item Un book de swaps dont l'échéance tombe sur un \omterm{def:rc:curve-construction:pillar}{pilier} aurait-il révélé quoi que ce soit~?
\end{enumerate}

\medskip\noindent\textbf{Partie III --- Les couvertures.}
\begin{enumerate}[resume]
\item Le bucket à quinze ans valait $-\text{USD}~346\,584$ par point de base et vaut désormais
$-583\,281$. Le swap au pair à quinze ans a un DV01 de 1\,149\,075~USD par milliard de notionnel.
Donner le notionnel de couverture avant et après.
\item Faire de même pour le bucket à vingt ans ($345\,841$ puis $465\,177$~; 1\,388\,775~USD
par milliard de notionnel).
\item La somme des buckets vaut 555\,270~USD avant et 557\,510 après. Pourquoi si
proches, alors que des buckets isolés ont bougé d'un tiers ou plus~?
\item Pourquoi les buckets ont-ils été calculés avec un bump d'un centième de point de base~?
\item Un nouveau bucket à douze ans apparaît ($-77\,741$). D'où vient-il~?
\end{enumerate}

\medskip\noindent\textbf{Partie IV --- Jugement.}
\begin{enumerate}[resume]
\item Donner deux arguments en faveur de la nouvelle interpolation et un contre.
\item Comment doit être encadré un changement de méthodologie de courbe~?
\item Quelle réserve un contrôleur pourrait-il constituer pour l'incertitude liée à la méthodologie de courbe~?
\item Énoncer le \emph{résultat nommé}~: le P\&L du passage à la nouvelle interpolation et la variation de la
couverture à quinze ans.
\item En une phrase~: que décide l'interpolation d'une courbe~?
\end{enumerate}
\end{problem}

\section{Questions d'entretien}

\begin{interviewq}[$\star$ \iqroles{chercheur, banque}]\label{iq:rc:curve-construction:1}
Quelle est la différence entre le bootstrap de courbe et un ajustement global de courbe, et
dans quels cas avez-vous besoin de la seconde~?
\end{interviewq}

\begin{interviewq}[$\star$ \iqroles{trader, chercheur}]\label{iq:rc:curve-construction:2}
Pourquoi l'interpolation linéaire des taux zéro donne-t-elle une mauvaise courbe à terme~?
\end{interviewq}

\begin{interviewq}[$\star\star$ \iqroles{chercheur, développeur}]\label{iq:rc:curve-construction:3}
Citez les propriétés que vous voudriez trouver dans une interpolation de courbe, et dites lesquelles
sont incompatibles.
\end{interviewq}

\begin{interviewq}[$\star\star$ \iqroles{trader}]\label{iq:rc:curve-construction:4}
Votre swap à huit ans montre un risque sur les buckets à quatre et à douze ans. Que
soupçonnez-vous~?
\end{interviewq}

\begin{interviewq}[$\star\star$ \iqroles{développeur, risque}]\label{iq:rc:curve-construction:5}
Comment gérez-vous le passage de fin d'année et les dates de réunion des banques centrales
dans un outil de construction de courbes~?
\end{interviewq}

\begin{interviewq}[$\star\star\star$ \iqroles{chercheur, risque}]\label{iq:rc:curve-construction:6}
Vos deltas par buckets changent lorsque vous modifiez la taille du bump, en passant d'un point de base
à un dixième de point de base. Le code est-il faux~?
\end{interviewq}
